% ===================================================================
%  اللوحة رقم ١٣ — الفندق. --- المطعم. المقهى.
%
%  Traduction, faite en 2026, en arabe levantin d'aujourd'hui, sur
%  l'ido de texto/io. Regles de la colonne : voir 10-lawha-01.tex.
%  Une seule numerotation. DEUX notes d'auteur, t13-noto-1 en tete et
%  t13-noto-2 posee entre les blocs 8 et 9. L'appel du bloc 9 se
%  compose en parentheses simples, « (*) », et non en exposant : c'est
%  ainsi dans l'ido, et le controle distingue les deux formes.
%
%  LE TROISIEME MOT DU TITRE EST UN MOT DE CETTE LANGUE. « La
%  Kafeerio » vient de café, qui vient de l'arabe قهوة, qui a donne
%  caffè, coffee, Kaffee, kофе et le reste. On ecrit « المقهى », de la
%  meme racine. La colonne gujaratie avait trouve la meme chose au
%  tableau 16 avec « બજાર » — le seul titre du livret que les deux
%  langues possedaient deja. Celle-ci l'a au tableau 13, et le mot
%  n'est pas passe par elle : IL EN EST SORTI.
%
%  ET LES ECHECS (78) SONT « الشطرنج », QUI EST LA ROUTE ELLE-MEME. Le
%  releve tient maintenant les deux bouts du meme mot : la colonne
%  gujaratie avait note que « શતરંજ » etait parti d'Inde — chaturanga,
%  les quatre corps d'une armee — vers l'Europe. Cette colonne-ci est
%  l'etape du milieu : le sanskrit est devenu شطرنج en arabe par le
%  persan, et c'est de cette forme-la que viennent scacchi, échecs,
%  chess. DEUX COLONNES DU RELEVE, DEUX BOUTS D'UNE SEULE ROUTE, ET
%  ELLES SE RENCONTRENT SUR UNE TABLE DE CAFE A ROYAN.
%
%  LE MENU DE L'ALINEA 9 PASSE ICI SANS UNE EGRATIGNURE, ET C'EST
%  PRECIEUX PARCE QUE C'EST LUI QUI AVAIT CASSE LA COLONNE GUJARATIE.
%  Crevettes, tripes, poisson, gigot d'agneau, epinards a la creme,
%  fromage, un Saint-Emilion : le Levant mange tout cela et le nomme
%  sans effort — قريدس, كرشة, سمك, فخد خروف, سبانخ, جبنة, نبيد. La
%  difficulte du tableau 13 gujarati n'etait donc PAS la francite du
%  repas, comme on aurait pu le croire : c'etaient les interdits. Une
%  table mediterraneenne se traduit d'un bout a l'autre de la
%  Mediterranee. IL AURA FALLU DEUX COLONNES POUR SAVOIR CE QUI ETAIT
%  DIFFICILE DANS LA PREMIERE.
%
%  ENFIN LE TRICTRAC (81) EST « الطاولة », QUI EST AUSSI LA TABLE. La
%  colonne egyptienne du releve a du inscrire ce renvoi-la dans ses
%  exemptions de kolonoj.py, parce que sa regle « طاولة -> ترابيزة »
%  criait sur le nom du jeu. Ici les deux sens sont justes tous les
%  deux et il n'y a rien a exempter : dans un cafe de Damas ou de
%  Beyrouth, on joue « طاولة » sur une طاولة, et personne ne s'y
%  trompe. UNE MEME AMBIGUITE EST UN PIEGE DANS UNE COLONNE ET UN
%  USAGE ORDINAIRE DANS LA VOISINE.
% ===================================================================

%%K t13-noto-1 noto
\VUnotes{143.2mm}{%
\textit{\VUgras{(*) لتفاصيل الغرفة، شوف اللوحة رقم ٦.}}%
}

%%K t13-apar-1 apar
\VUsaut{3.0mm}
\VUtitre{15.0pt}{60}{اللوحة رقم ١٣}
\VUsaut{1.6mm}
\VUpk{10.2pt}{40}{الفندق. --- المطعم.}
\VUsaut{2.0mm}
\VUpk{10.2pt}{40}{المقهى.}
\VUsaut{3.4mm}

%%K t13-01-1 p
1. --- وصلت مبارح عالمسا على رويان، ونزلت بـ« \textit{فندق المحيط} »
اللي نصحني فيه رفيق.

%%K t13-01-2 p
أعطوني \VUgras{غرفة} \textsuperscript{(2)} حلوة، عالـ\VUgras{طابق}
\textsuperscript{(1)} التاني، ونمت فيها كتير منيح (*).

%%K t13-02-1 p
2. --- واليوم الصبح، وأنا طالع من غرفتي، اللي جابت عليها
\VUgras{الصبيّة} \textsuperscript{(3)} الفطور، فتّ على \VUgras{غرفة
التدخين} \textsuperscript{(5)} اللي بتنستعمل كمان \VUgras{غرفة مطالعة}
\textsuperscript{(4)} لقراية \VUgras{الجرايد} \textsuperscript{(6)}، وأنا
عم دخّن \VUgras{سيجارة} \textsuperscript{(8)}. وبعد ما قريت، نزلت
بـ\VUgras{الأصانسير} \textsuperscript{(9)} ورحت اتمشّى بالمدينة.

%%K t13-03-1 p
3. --- ولأني نسيت اسأل \VUgras{موظّف} \textsuperscript{(12)} الفندق أي
ساعة بيقدّموا الأكل عالـ\VUgras{سفرة المشتركة} \textsuperscript{(11)}،
رجعت بكّير واستفدت لروح استحمّي بـ\VUgras{حمّام} \textsuperscript{(10)}
الفندق. وبعدين رجعت عالـ\VUgras{صندوق} \textsuperscript{(15)} لتلفن
لواحد من رفقاتي. بس \VUgras{التلفون} \textsuperscript{(13)} كان مشغول ؛
ولمّا شافت حيرتي \VUgras{الصرّافة} \textsuperscript{(14)}، اللي كانت عم
تسجّل \VUgras{فاتورة} \textsuperscript{(17)} على \VUgras{سجلّ النزلا}
\textsuperscript{(16)}، طلبت من \VUgras{البوّاب} \textsuperscript{(18)}
يبعت \VUgras{الصبي} \textsuperscript{(19)} يعمل شغلتي
\VUgras{عالبسكليت} \textsuperscript{(20)}.

%%K t13-04-1 p
4. --- وشكرتها وطلعت لأعطي بقشيش لـ\VUgras{الحمّال}
\textsuperscript{(24)} (العتّال) اللي جاب من المحطّة على
\VUgras{عربايته} \textsuperscript{(25)} \VUgras{علبة برنيطتي}
\textsuperscript{(26)}، و\VUgras{شنتتي} \textsuperscript{(27)}
و\VUgras{كيس سفري} \textsuperscript{(28)}. و\VUgras{الساعي}
\textsuperscript{(21)} اللي هلّق وصل، أعطاني \VUgras{مكتوب}
\textsuperscript{(22)}.

%%K t13-05-1 p
5. --- وضلّيت كمان شوي على عتبة الباب، جنب \VUgras{صندوق البريد}
\textsuperscript{(23)}، عم اتفرّج عالحركة بالشارع. و\VUgras{سائق}
\textsuperscript{(30)} \VUgras{الأومنيبوس} \textsuperscript{(29)} كان عم
يحمّل على عربايته \VUgras{صندوق} \textsuperscript{(31)}
\VUgras{مسافر} \textsuperscript{(32)} كان عم يودّعه \VUgras{صاحب
الفندق} \textsuperscript{(33)}. و\VUgras{موزّع برقيات}
\textsuperscript{(35)} شاب جاب على مهله \VUgras{برقية}
\textsuperscript{(34)} لزبون من زباين الفندق ؛ و\VUgras{سايح}
\textsuperscript{(36)}، و\VUgras{الكاميرا} \textsuperscript{(38)} معلّقة
عكتفه، كان عم يتطلّع بـ\VUgras{دليله} \textsuperscript{(37)}.

%%K t13-tit-1 sub
\VUsaut{4.0mm}
\VUpk{10.2pt}{40}{وقت الغدا.}
\VUsaut{3.0mm}

%%K t13-06-1 p
6. --- وقدّام السياج، \VUgras{مشوارجي} \textsuperscript{(39)} ما
بيهتمّ كان عم يغفى بين \VUgras{كلّاب حمله} \textsuperscript{(40)}
و\VUgras{علبة البويا} \textsuperscript{(41)} تبعه، وقت \VUgras{غسّالة}
\textsuperscript{(42)} شابّة كانت شايلة \VUgras{سلّة} \textsuperscript{(43)}
بياضات، عم تضحك من الوجّ الجدّي تبع \VUgras{رئيس الصالة}
\textsuperscript{(44)} اللي كان واقف جنب الباب.

%%K t13-07-1 p
7. --- ولمّا شفت \VUgras{الطبّاخ} \textsuperscript{(45)} طالع من
\VUgras{مطبخه} \textsuperscript{(47)} اللي عالـ\VUgras{طابق التحتاني}
\textsuperscript{(48)} ليبعت \VUgras{صبي المطبخ} \textsuperscript{(46)}
بشغلة، تذكّرت إنّه وقت الغدا قرّب وفتّ عالمطعم، وأنا عم اقطع الحوش
اللي فيه \VUgras{سائق} \textsuperscript{(50)} كان مركّن \VUgras{سيّارته}
\textsuperscript{(49)}، وقت \VUgras{ميكانيكي} \textsuperscript{(52)} كان
عم يصلّح، حسب تعليمات \VUgras{صاحب البسكليت} \textsuperscript{(53)}،
\VUgras{ثلاثية عجلات} \textsuperscript{(51)} تعمل بالبنزين وهلّق صار فيها
حادث.

%%K t13-08-1 p
8. --- وسألت \VUgras{صبي} \textsuperscript{(54)} شاب كان شايل
\VUgras{كاسات بيرة} \textsuperscript{(55)}، إذا السفرة المشتركة تحت
الفرندة، ولمّا جاوبني إي، قعدت على وحدة من الطاولات وخلّيتهن يقدّموا
لي أكل ممتاز.

%%K t13-noto-2 noto
\VUnotes{143.2mm}{%
\textit{\VUgras{(*) بمساعدة قائمة الأطباق الموجودة بالقاموس، بيقدر
المعلّم يغيّر بسهولة بالمنيو اللي تحت.}}%
}

%%K t13-tit-2 sub
\VUsaut{4.0mm}
\VUpk{10.2pt}{40}{غدا ممتاز.}
\VUsaut{3.0mm}

%%K t13-09-1 p
9. --- الجرسون جابلي منيو (*) الغدا وقائمة النبيد. واخترت وطلبت
كـ\VUgras{مقبّلات قريدس} مع \VUgras{زبدة}، وكـ\VUgras{طبق أول}
\VUgras{كرشة على طريقة كان}. وما أكلت \VUgras{سمك}، بس طلبت قطعة
\VUgras{فخد} خروف، وكـ\VUgras{طبق خضرة} \VUgras{سبانخ} مع
\VUgras{قشطة}. وبعدين، لأنه ما عجبني شي من \VUgras{الأطباق
الوسطانية}، خلّيتهن يجيبولي حلويات مختلفة — \VUgras{جبنة}،
\VUgras{فواكه} و\VUgras{بسكوت}. وكمان زدت \VUgras{نبيد} سانت-إميليون
ممتاز، وأنا كتير راضي عن أكلي تركت الطاولة بعد ما أعطيت
\VUgras{الجرسون} \textsuperscript{(57)} \VUgras{بقشيش}
\textsuperscript{(59)} منيح.

%%K t13-10-1 p
10. --- ولمّا شافني \VUgras{المدير} \textsuperscript{(60)} جاهز
للطلوع، اقترح عليّ روح عالشرفة لأتأمّل منظر \VUgras{المحيط}
\textsuperscript{(62)} الرائع. وبعيد بلّشت تبيّن \VUgras{قوارب صيد}
\textsuperscript{(63)}، و\VUgras{مراكب شراعية} \textsuperscript{(64)}،
و\VUgras{باخرة ركّاب} \textsuperscript{(65)} ضخمة، اللي ظلّها الأسود عم
يبرز عالـ\VUgras{غيوم} \textsuperscript{(67)} البيضا تبع \VUgras{الأفق}
\textsuperscript{(66)}. ونسمة خفيفة كانت عم تهزّ \VUgras{الراية}
\textsuperscript{(70)} المركّزة فوق \VUgras{لافتة} \textsuperscript{(69)}
\VUgras{القاعة} \textsuperscript{(68)}.

%%K t13-tit-3 sub
\VUsaut{3.0mm}
\VUpk{10.2pt}{40}{التسلية.}
\VUsaut{3.0mm}

%%K t13-11-1 p
11. --- والمنظر كان كتير مشوّق لدرجة إنّي قرّرت اطلع بكرا على سطح
المقهى وآخد معي \VUgras{منظار مزدوج} \textsuperscript{(71)} أو
\VUgras{منظار} \textsuperscript{(72)}.

%%K t13-12-1 p
12. --- ولمّا تركت الشرفة، رحت على مقهى الفندق لشرب \VUgras{مشروب}
\textsuperscript{(73)} وألعب \VUgras{دورة بلياردو} \textsuperscript{(75)}.
ومرقت بلا ما وقف بصالة القهوة اللي فيها زباين كتار عم يلعبوا
\VUgras{شطرنج} \textsuperscript{(78)}، و\VUgras{ضاما} \textsuperscript{(79)}،
و\VUgras{دومينو} \textsuperscript{(80)}، و\VUgras{طاولة}
\textsuperscript{(81)} أو \VUgras{شدّة} \textsuperscript{(82)}، وطلعت
دغري عالـ\VUgras{صالة البلياردو} \textsuperscript{(74)}.

%%K t13-13-1 p
13. --- ولمّا خلصت الدورة، نزلت عالطابق الأرضي وطلبت من الجرسون
\VUgras{ظرف} \textsuperscript{(84)}، و\VUgras{ورق} \textsuperscript{(83)}،
و\VUgras{طابع بريد} \textsuperscript{(85)}، و\VUgras{قلم}
\textsuperscript{(87)} و\VUgras{حبر} \textsuperscript{(86)} لكتب مكتوب.

%%K t13-13-2 p
وبعدين ناديت \VUgras{عربجي} \textsuperscript{(92)} \VUgras{عربايته}
\textsuperscript{(88)} كانت واقفة قدّام المقهى. وسألته عن السعر بالساعة
أو بالجولة، ولمّا اتّفقنا عالسعر، فتحلي \VUgras{باب العربايّة}
\textsuperscript{(89)} وقعّدني. ورجع طلع عالـ\VUgras{مقعد}
\textsuperscript{(91)}، وحرّك الخيل بسرعة بـ\VUgras{الكرباج}
\textsuperscript{(93)}، وسافرنا بلفّة ساحرة.

\VUsaut{6.0mm}
\VUfilet{22mm}
